% =============================================================================
% unibs-style.tex -- UNIBS visual identity for the document body.
% Included through `include-in-header` by both pandoc defaults files.
% Sources: res/reference/UNIBS_Branding-Guidelines_extract.pdf
%   - palette primaria (Blu notte, Blu scuro, Neutro caldo, Bianco, Nero + tints)
%   - font: Antarctica, fallback Helvetica Neue (body 10.5 pt, ragged right)
%   - visual elements: orthogonal divider lines, rounded joints
% Students never edit this file.
% =============================================================================
\makeatletter

% --- Colours -----------------------------------------------------------------
\definecolor{unibsBlue}{HTML}{27326E}      % Blu notte
\definecolor{unibsDarkBlue}{HTML}{242C59}  % Blu scuro
\definecolor{unibsNeutral}{HTML}{F7F5EB}   % Neutro caldo
\definecolor{unibsBlack}{HTML}{1D1D1B}     % Nero
% Blu notte percentage tints, sampled from the brand palette (not white mixes).
\definecolor{unibsBlue5}{HTML}{E2DFDA}
\definecolor{unibsBlue10}{HTML}{D6D4D4}
\definecolor{unibsBlue20}{HTML}{BEBCC4}
\definecolor{unibsBlue30}{HTML}{A7A5B6}
\definecolor{unibsBlue40}{HTML}{9392AA}
\definecolor{unibsBlue50}{HTML}{81809F}
\definecolor{unibsBlue60}{HTML}{6F6E94}
\definecolor{unibsBlue70}{HTML}{5E608A}
\definecolor{unibsBlue80}{HTML}{4D5281}
\definecolor{unibsBlue90}{HTML}{3C4579}

% --- Fonts ---------------------------------------------------------------------
% Antarctica (brand font, commercial) when installed; otherwise Helvetica Neue
% (closest system match, ships with macOS); otherwise TeX Gyre Heros
% (Helvetica clone shipped with TeX Live/MiKTeX, used e.g. on Windows).
% Antarctica has no italic: a slanted upright is used, as for UNIBS documents.
%
% Besides regular (series m) and bold (b), every installed weight of the family
% is registered for the cover tuning (\unibsCoverStyle in unibs-covers.tex):
% "<family> <weight>" -> NFSS series, e.g. "Helvetica Neue Light" -> l.
\RequirePackage{fontspec}
\def\unibs@faces{}
\newcommand\unibs@face[3]{% #1 family, #2 NFSS series, #3 weight name
  \IfFontExistsTF{#1 #3}
    {\edef\unibs@faces{\unibs@faces FontFace={#2}{n}{Font={#1 #3}},}}{}}
\newcommand\unibs@setfonts[2]{% #1 family, #2 extra fontspec options
  \unibs@face{#1}{ul}{UltraLight}\unibs@face{#1}{el}{Thin}%
  \unibs@face{#1}{l}{Light}\unibs@face{#1}{sb}{Medium}%
  \unibs@face{#1}{ub}{Black}%
  \edef\unibs@next{%
    \noexpand\setmainfont{#1}[#2\unibs@faces]%
    \noexpand\setsansfont{#1}[#2\unibs@faces]}%
  \unibs@next}
\newif\ifunibs@antarctica
\IfFontExistsTF{Antarctica}{\unibs@antarcticatrue}{}
\ifunibs@antarctica
  \IfFontExistsTF{Antarctica Bold}
    {\def\unibs@bold{Antarctica Bold}}
    {\def\unibs@bold{Antarctica}}
  \unibs@setfonts{Antarctica}{%
    BoldFont=\unibs@bold,
    ItalicFont=Antarctica, ItalicFeatures={FakeSlant=0.18},
    BoldItalicFont=\unibs@bold, BoldItalicFeatures={FakeSlant=0.18},}
\else
  \IfFontExistsTF{Helvetica Neue}
    {\unibs@setfonts{Helvetica Neue}{}}
    {\unibs@setfonts{TeX Gyre Heros}{}%
     \PackageWarningNoLine{unibs}{Neither Antarctica nor Helvetica Neue found:
       using TeX Gyre Heros}}
\fi
\IfFontExistsTF{DejaVu Sans Mono}
  {\setmonofont{DejaVu Sans Mono}[Scale=MatchLowercase]}
  {\setmonofont{Latin Modern Mono}[Scale=MatchLowercase]}

% --- Page layout, header and footer -------------------------------------------
\RequirePackage{etoolbox}
\RequirePackage{ragged2e}
\RequirePackage{scrlayer-scrpage}
\providecommand\unibsHeadLeft{}
\providecommand\unibsStudentName{}
\clearpairofpagestyles
\ihead{\unibsHeadLeft}
\ohead{\unibsStudentName}
\ifoot{\textbf{\textcolor{unibsBlue}{Universit\`a degli Studi di Brescia}}}
\ofoot{\pagemark}
\setkomafont{pageheadfoot}{\footnotesize\color{unibsBlue60}}
\setkomafont{pagenumber}{\footnotesize\bfseries\color{unibsBlue}}
\KOMAoptions{headsepline=0.4pt, footsepline=0.4pt}
\setkomafont{headsepline}{\color{unibsBlue10}}
\setkomafont{footsepline}{\color{unibsBlue10}}
\pagestyle{scrheadings}

% Brand manual: body text flush left / ragged right ("bandiera sinistra").
\AtBeginDocument{\RaggedRight\color{unibsBlack}}

% --- Headings --------------------------------------------------------------------
\setkomafont{disposition}{\bfseries\color{unibsBlue}}
\RedeclareSectionCommand[
  font=\fontsize{17}{21}\selectfont,
  beforeskip=2.4\baselineskip plus .4\baselineskip,
  afterskip=.9\baselineskip]{section}
\RedeclareSectionCommand[
  font=\fontsize{13}{16}\selectfont,
  beforeskip=1.6\baselineskip plus .3\baselineskip,
  afterskip=.5\baselineskip]{subsection}
\RedeclareSectionCommand[
  font=\fontsize{10.5}{13}\selectfont,
  beforeskip=1.2\baselineskip plus .2\baselineskip,
  afterskip=.3\baselineskip]{subsubsection}
\renewcommand*\sectionformat{\textcolor{unibsBlue40}{\thesection}\hskip .7em}
\renewcommand*\subsectionformat{\textcolor{unibsBlue40}{\thesubsection}\hskip .6em}
\renewcommand*\subsubsectionformat{\textcolor{unibsBlue40}{\thesubsubsection}\hskip .5em}
% Level-1 headings open with a divider line (brand "linee divisorie").
% \unibsNoSectionLine drops it for the next section only: the compendium
% filter emits it when a report starts with a section right after its cover,
% where the line would separate nothing.
\newif\ifunibs@nosecline
\newcommand\unibsNoSectionLine{\global\unibs@noseclinetrue}
\renewcommand*\sectionlinesformat[4]{%
  \ifstrequal{#1}{section}{%
    \ifunibs@nosecline
      \global\unibs@noseclinefalse
    \else
      {\color{unibsBlue10}\rule[.5\baselineskip]{\linewidth}{.5pt}}\par\nobreak
      \vspace{-.15\baselineskip}%
    \fi}{}%
  \@hangfrom{\hskip #2#3}{#4}}

% --- Lists ------------------------------------------------------------------------
\RequirePackage{enumitem}
\setlist{itemsep=.25em, topsep=.4em}
\setlist[itemize,1]{label=\textcolor{unibsBlue}{\small\textbullet}}
\setlist[itemize,2]{label=\textcolor{unibsBlue40}{\textendash}}
\setlist[itemize,3]{label=\textcolor{unibsBlue40}{\small\textbullet}}
\setlist[enumerate,1]{label=\textcolor{unibsBlue}{\bfseries\arabic*.}}
\setlist[enumerate,2]{label=\textcolor{unibsBlue}{\alph*)}}

% --- Captions, figures, tables ------------------------------------------------
\RequirePackage{caption}
\captionsetup{font={small}, labelfont={bf,color=unibsBlue}, textfont={color=unibsBlue80},
  labelsep=period, justification=raggedright, singlelinecheck=true}
\captionsetup[table]{skip=6pt}
\RequirePackage{booktabs}
\RequirePackage{longtable}
\RequirePackage{colortbl}
\AtBeginDocument{\arrayrulecolor{unibsBlue}}
\setlength\heavyrulewidth{.8pt}
\setlength\lightrulewidth{.4pt}
\AtBeginEnvironment{longtable}{\small}
% Header cells: the Lua filter inserts \unibsTableHead at the start of each cell.
\newcommand\unibsTableHead{\bfseries\color{unibsBlue}}

% --- Quotes and code ----------------------------------------------------------------
\RequirePackage{tcolorbox}
\tcbuselibrary{skins,breakable}
\renewenvironment{quote}
  {\begin{tcolorbox}[blanker, breakable, left=5mm, top=1mm, bottom=1mm,
     borderline west={2pt}{0pt}{unibsBlue20}, before skip=.6\baselineskip,
     after skip=.6\baselineskip]\color{unibsBlue80}}
  {\end{tcolorbox}}
% Every code block is wrapped in `unibscode` by the Lua filter.
\newtcolorbox{unibscode}{enhanced, breakable, colback=unibsNeutral, colframe=unibsNeutral,
  boxrule=0pt, arc=1.5mm, left=3mm, right=3mm, top=1.5mm, bottom=1.5mm,
  borderline west={2pt}{0pt}{unibsBlue}, fontupper=\small,
  before skip=.6\baselineskip, after skip=.8\baselineskip}
\AtBeginDocument{%
  \@ifundefined{Shaded}{}{\renewenvironment{Shaded}{}{}}%
  \let\unibs@texttt\texttt
  \renewcommand\texttt[1]{\textcolor{unibsBlue}{\unibs@texttt{#1}}}}

% --- Misc ---------------------------------------------------------------------
\renewcommand\footnoterule{{\color{unibsBlue10}\rule{.3\linewidth}{.4pt}}\vspace{2pt}}
\deffootnote{1.2em}{1em}{\textcolor{unibsBlue}{\thefootnotemark}\hskip .4em}
\AtBeginDocument{\urlstyle{same}}

\makeatother
